\documentclass[a4paper, 12pt]{article}
\usepackage{graphicx}
\usepackage[T2A]{fontenc} 
\usepackage[utf8]{inputenc}
\usepackage[english, russian]{babel}
\usepackage{amsmath}
\usepackage{xcolor}
\usepackage{array}
\usepackage[left=2cm,right=2cm,top=1cm,bottom=0.5cm,bindingoffset=0cm]{geometry}
\usepackage{float}
\usepackage{wrapfig}
\usepackage{tabularx}
\usepackage{nicematrix}
\usepackage{multirow}
\usepackage{multicol}

\begin{document}

\section{Билет 4. Гомоморфизмы колец. Ядро и образ гомоморфизма}
Гомоморфизм (единая форма) - это функция, при которой у каждого аргумента есть только один образ и при этом выполняется:\\

\begin{center}
	$f(a + b) = f(a) + f(b)$ и $f(a \cdot b) = f(a) \cdot f(b)$
\end{center}

Очевидным примером будет $f(x) = x, f: Z \rightarrow Z $, где все свойства будут выполняться. \\
Функции вида $f(x) = ax$ не будут подходить по умножению\\

Ядро (корни) гомоморфизма f — это $Ker(f ) = \{ x \in K : f(x) = 0 \}$.\\
Ядро гомоморфизма - это такое множество, которое содержит все элементы кольца, при котором отображение этого элемнта будет нулем.\\

Если гомоморфизм $f(x)$, то $Ker(f) = {0}$

Образ гомоморфизма - то же самое, что и образ функции: множество всех отображений.

$Im(f) = Z$

Пусть $K, L$ — кольца, $f : K \rightarrow L$ — гомоморфизм колец.\\
Тогда:\\
1) $Ker(f )$ — подкольцо K.
2) $Im(f )$ — подкольцо L.\\
Доказательство.\\
1) Пусть $a, b \in Ker(f)$. Тогда\\
а) $f (a + b) = f (a) + f (b) = 0 + 0 = 0$, следовательно,\\
$a + b \in Ker(f )$.
\\
б) $f (a \cdot b) = f(a) \cdot f(b) = 0 \cdot 0 = 0$, следовательно,
$a \cdot b \in Ker(f )$.
\\
в) $f (-a) = -f (a) = -0_{L} = 0_{L}$.\\

2) Пусть $y_{1}, y_{2} \in Im(f)$, а $x_{1}, x_{2} \in K$ таковы, что $f(x_{1}) = y_2$ и $f(x_2) = y_2$.
Тогда \\
а) $y_1 + y_2 = f(x_1) + f(x_2) = f (x_1 + x_2) \in Im(f)$ и $y_1 \cdot y_1 = f(x_1) \cdot f(x_2) \in Im(f).$
\\
б) $-y = -f(x) = f (-x) \in Im(f)$.

\section{Билет 5. Типы гомоморфизмов. Мономорфизм и ядро}

Типы гомоморфизма (как в матане)
\\
Инъекция - мономорфзима\\
Суръекция - эпиморфизм\\
Биекция - изоморфизм\\

Инъекция, суръекция, биекция - для описания функций.\\
мономорфизм, эпиморфизм, изоморфизм - для описания отображений.


\end{document}
